% =====================================================================
% ViAmpleHate — ACL paper (VERSION v_review: v3 + revision rounds 1-2)
% Revision applied per review_report.md (root). Round 2 filled real numbers from the
% notebooks (hyperparameters, baseline accuracies, cue-bank sizes, target/attack
% coverage, VOZ per-class) and removed the ablation section per author instruction.
% Remaining open items (single-run only / no significance test; no published-SOTA
% comparison) are stated in Limitations. NO numeric results were fabricated;
% v3 is left untouched. See REVISION_NOTES.md.
% Drop into the Overleaf ACL Conference template, replacing acl_latex.tex.
% Requires: acl.sty, acl_natbib.bst (template) + custom.bib (this folder). pdfLaTeX.
% Compiles with pdfLaTeX (author names are ASCII; cue examples are ASCII-safe).
% For Vietnamese text with diacritics, use XeLaTeX/LuaLaTeX or the T5 fontenc line below.
% =====================================================================
\documentclass[11pt]{article}

\usepackage[final]{acl}  % shows authors (use "review" for anonymized submission)

\usepackage{times}
\usepackage{latexsym}
\usepackage[T1]{fontenc}
% \usepackage[T5]{fontenc} % Uncomment (and drop T1) for full Vietnamese coverage
\usepackage[utf8]{inputenc}
\usepackage{microtype}
\usepackage{inconsolata}
\usepackage{graphicx}
\usepackage{amsmath}
\usepackage{amssymb}
\usepackage{booktabs}
\usepackage{array}
\usepackage{tikz}
\usetikzlibrary{arrows.meta,positioning,fit,backgrounds,calc}

\title{ViAmpleHate: A Proposed AmpleHate- and PhoBERT-based Approach\\ for Vietnamese Hate Speech Detection}

\author{
  \textbf{Tran Nguyen Duc Trung} \quad
  \textbf{Nguyen Ha Minh Tuan} \quad
  \textbf{Nguyen Gia Cat Long} \\
  \textbf{Tran Phan Thanh Tung} \quad
  \textbf{Mo Van Tung} \\[4pt]
  University of Information Technology, VNU-HCM, Ho Chi Minh City, Vietnam \\
  \texttt{\{23521687, 23521718, 23520881, 23521747, 23521741\}@gm.uit.edu.vn}
  % Supervisor: Huynh Van Tin (last author) -- TODO: add faculty email (e.g. ...@uit.edu.vn)
}

\setlength\titlebox{6cm}

\begin{document}
\maketitle

\begin{abstract}
Detecting hate speech on Vietnamese social media is difficult because hateful intent is rarely carried by explicit named entities; instead it hides in informal group references, slang, sarcasm, and indirect insinuation. Target-aware models such as AmpleHate have shown that attending to the relationship between an utterance and the target it addresses helps detect even implicit hate, yet AmpleHate's original pipeline relies on English named-entity recognition (NER) and English-oriented target categories, so it transfers poorly to Vietnamese: on our training data English NER finds a usable target in only about 0.09\% of comments, collapsing the model into a plain sentence classifier. We propose \textbf{ViAmpleHate}, a Vietnamese adaptation of AmpleHate built on PhoBERT. ViAmpleHate (i) replaces English NER with Vietnamese NER and a curated \emph{target-cue} bank, raising target coverage from $\approx$0.09\% to $\approx$19--45\% depending on the dataset; (ii) adds a separate \emph{attack-cue} bank and models three relation channels---explicit target, implicit context, and attack---through a \emph{relation-bank attention}; (iii) replaces fixed-scalar injection with an \emph{instance-adaptive gate}; and (iv) trains with weighted cross-entropy combined with a contrastive objective. On ViHSD and VOZ-HSD in a binary NON-HATE/HATE setting, ViAmpleHate improves macro-F1 and the minority HATE-class F1 over an AmpleHate baseline and over TF-IDF, BiLSTM, and PhoBERT-CNN baselines, with the largest and clearest gains on the minority class (HATE-F1 $+0.0508$ on VOZ-HSD); improvements on ViHSD are small and we report them as preliminary, pending a controlled, multi-seed evaluation. We further analyse how target-cue coverage relates to these gains, supporting---rather than proving---the view that target-aware modeling should be adapted to Vietnamese rather than transferred verbatim from English.
\end{abstract}

\section{Introduction}

The rapid growth of Vietnamese-language social media has been accompanied by a corresponding spread of hateful and abusive content. Such content harms targeted individuals and communities, degrades online discourse, and creates legal and reputational risk for platforms. Because the volume of user-generated text vastly exceeds what human moderators can review, automatic hate speech detection has become a practical necessity. Yet most progress has centred on English and other high-resource languages, and methods that work well there often degrade when applied to Vietnamese.

Vietnamese hate speech detection is hard for three intertwined reasons. First, social-media language is extremely noisy: users abbreviate, write in \emph{teencode}, elongate characters for emphasis, embed emojis with pragmatic meaning, and---deliberately or not---misspell words, including obfuscating profanity to evade filters. Second, and most importantly for this work, the \emph{target} of an attack is rarely a named entity. Whereas English hate speech frequently names a person, organization, or nationality, Vietnamese hostility is typically aimed at a group through pronouns, kinship or group nouns, regional labels, or informal forms of address. Third, the data is severely imbalanced: in both datasets we study, the HATE class accounts for only about one tenth of all comments, so a model can reach high accuracy while detecting almost no hate.

A promising line of work models hate speech in a \emph{target-aware} manner. AmpleHate \citep{amplehate}, in particular, shows that explicitly attending to the relationship between a sentence and its potential target improves detection of \emph{implicit} hate, where hostility is conveyed without overt slurs. AmpleHate extracts candidate targets with NER, computes a target-aware attention vector, and injects it into the sentence representation before classification. This is attractive for Vietnamese, where the target is so central---but the original pipeline is built around English NER and English target categories. When we run it on Vietnamese with an English NER component, it finds a usable target in only 21 of 24{,}048 training comments (about 0.09\%). For the overwhelming majority of inputs the model falls back to the global \texttt{[CLS]} representation and behaves like an ordinary PhoBERT classifier, discarding precisely the target-aware signal that motivated it.

We propose \textbf{ViAmpleHate}, a Vietnamese adaptation of AmpleHate that re-engineers each stage of the pipeline around how Vietnamese speakers actually express hate. Our contributions are:
\begin{enumerate}
  \item \textbf{Vietnamese target extraction.} We replace English NER with a Vietnamese NER model and a curated \emph{target-cue} bank of pronouns, group nouns, and informal address forms, raising the fraction of comments with a detected target from $\approx$0.09\% to $\approx$18.8--20.0\%.
  \item \textbf{Separate target and attack channels.} We add an \emph{attack-cue} bank for hostile predicates and model target and attack cues as distinct signals through a three-channel \textbf{relation-bank attention}.
  \item \textbf{Adaptive fusion (plus an implementation note).} We replace AmpleHate's fixed-scalar injection with an \textbf{instance-adaptive gate}. We additionally describe an implementation-level correction to a batch-level attention computation in our port, which we report as an implementation note rather than a primary contribution.
  \item \textbf{Training objective.} We combine weighted cross-entropy with a contrastive loss that sharpens HATE/NON-HATE separation under imbalance.
  \item \textbf{Empirical study.} We evaluate on ViHSD and VOZ-HSD against five baselines and analyse where and why the gains arise, including a target-coverage analysis and a qualitative study of remaining errors.
\end{enumerate}
The broader message is that target-awareness is a genuinely useful inductive bias for hate speech detection, but only when adapted to the linguistic surface of the target language instead of being transferred verbatim from English.

\section{Related Work}

\subsection{Hate speech detection and implicit hate}
Automatic hate speech detection has evolved from lexicon- and feature-based classifiers to deep sequence models and, more recently, pretrained transformers. A persistent difficulty is \emph{implicit} hate: messages that convey hostility without explicit slurs, relying on sarcasm, stereotypes, coded references, or in-group knowledge. Implicit hate is often more common than overt slurs and substantially harder for keyword-driven systems, motivating approaches that reason about \emph{who} is targeted and \emph{how}.

\subsection{AmpleHate}
AmpleHate \citep{amplehate} addresses implicit hate by amplifying the attention between a sentence's global representation and its candidate targets. It extracts target tokens with NER, computes a HeadAttention vector whose query derives from \texttt{[CLS]} and whose keys/values derive from target tokens, and injects the result with a fixed coefficient before classification. Its central insight---that the target is a load-bearing signal---directly motivates our work. Its limitations for Vietnamese are equally direct: it presumes English-style NER and target categories, injects relation information at a fixed strength for every instance, and models a single relation (target) while ignoring the hostile predicate, or \emph{attack}, that turns a mention into hate.

\subsection{Vietnamese hate speech detection}
Vietnamese has seen growing interest in hate speech resources and models. ViHSD \citep{vihsd} provides a large-scale three-way (CLEAN/OFFENSIVE/HATE) dataset and is a standard benchmark. ViTHSD \citep{vithsd} reframes the task around \emph{targets}, annotating the degree of hatred directed at each target; this reinforces the centrality of the target in Vietnamese hate speech and supports our target-aware design. Unlike ViTHSD, ViAmpleHate requires no span-level target annotation; it locates target and attack positions with NER plus cue banks, which keeps it applicable to datasets carrying only sentence-level labels.

\subsection{PhoBERT and Vietnamese pretrained models}
PhoBERT \citep{phobert} is a RoBERTa-style model pretrained on large Vietnamese corpora and is the de facto encoder for Vietnamese text classification. Because PhoBERT is trained on \emph{word-segmented} text, word segmentation precedes tokenization. Hybrid architectures such as PhoBERT-CNN add local feature extractors on top of contextual embeddings. We use PhoBERT-base as the shared encoder for both baseline and proposed models so that differences are attributable to target-aware modeling rather than the backbone.

\subsection{Contrastive learning for text classification}
Contrastive objectives encourage same-class representations to be close and different-class representations to be far apart, improving robustness and class separation, particularly under imbalance \citep{supcon}. We adopt a contrastive auxiliary objective on the post-gate representation alongside weighted cross-entropy---a pairwise cosine-margin term inspired by supervised contrastive learning, not the exact SupCon loss---to tighten the boundary between the frequent NON-HATE class and the rare HATE class.

\section{Datasets}

\subsection{ViHSD}
ViHSD \citep{vihsd} is a Vietnamese hate speech dataset of social-media comments annotated with three original labels (CLEAN, OFFENSIVE, HATE) and pre-split into train/validation/test. After the binary relabeling of Section~\ref{sec:reshape}, the per-split distribution is in Table~\ref{tab:data}.

\subsection{VOZ-HSD}
VOZ-HSD \citep{vozhsd}, released with ViHateT5 \citep{vihatet5}, consists of comments from the VOZ forum (sourced from the public \texttt{tarudesu/VOZ-HSD} release). As with ViHSD, the data is split into train/validation/test, and its post-relabeling distribution is in Table~\ref{tab:data}. As a third-party resource, its detailed annotation procedure and licensing are documented by the dataset authors.

\subsection{Binary relabeling and preprocessing}
\label{sec:reshape}
The only modification we make is \emph{relabeling}; we do not alter or subsample the number of examples. We cast the task as binary classification by merging CLEAN and OFFENSIVE into NON-HATE and keeping HATE as the positive class. This focuses the problem on group-directed hate, as opposed to profanity or aggression in general, and yields a consistent label space across datasets. Each comment is then normalized (lowercasing; URL/symbol removal; collapsing elongated runs; teencode normalization; mapping frequent emojis to coarse pragmatic tags such as mockery, anger, disgust, or laughter), word-segmented---mandatory because PhoBERT is pretrained on word-segmented Vietnamese---and tokenized with the PhoBERT tokenizer.

\begin{table}[t]
  \centering
  \small
  \begin{tabular}{llrrr}
    \toprule
    \textbf{Dataset} & \textbf{Split} & \textbf{NON-HATE} & \textbf{HATE} & \textbf{Total} \\
    \midrule
    ViHSD & Train & 21{,}492 & 2{,}556 & 24{,}048 \\
    ViHSD & Dev   & 2{,}402  & 270     & 2{,}672  \\
    ViHSD & Test  & 5{,}992  & 688     & 6{,}680  \\
    VOZ   & Train & 26{,}993 & 3{,}007 & 30{,}000 \\
    VOZ   & Dev   & 4{,}520  & 480     & 5{,}000  \\
    VOZ   & Test  & 4{,}487  & 513     & 5{,}000  \\
    \bottomrule
  \end{tabular}
  \caption{Dataset statistics after binary relabeling. HATE is $\approx$10\% throughout, so macro-F1/HATE-F1 matter more than accuracy.}
  \label{tab:data}
\end{table}

% FIGURE 1 — Class distribution. DRAW NEW (matplotlib) from Table 1.
\begin{figure}[t]
  \centering
  \includegraphics[width=\columnwidth]{example-image-golden}
  \caption{Class distribution of ViHSD and VOZ-HSD after binary relabeling. \textbf{[DRAW NEW]}}
  \label{fig:dist}
\end{figure}

\subsection{Cue bank construction}
\label{sec:cue}
The target-cue and attack-cue banks are central to ViAmpleHate. Both are seeded manually from inspection of the training split and from common Vietnamese referential and offensive expressions, then normalized to match the preprocessing pipeline. The \textbf{target-cue bank} contains referential expressions that often introduce a person or group---pronouns, kinship and group nouns, regional and demographic labels, and informal address forms---but that are \emph{not} by themselves indicators of hate. The \textbf{attack-cue bank} contains hostile predicates: insults, dehumanizing terms, threats, and strongly negative evaluations. Keeping the banks separate is deliberate: a target mention without a hostile predicate is usually not hate, and a hostile predicate without a clear group target is often mere offensiveness. Each cue is normalized, word-segmented, and tokenized, then matched against the token stream by full token-sequence matching so a multi-token Vietnamese phrase is matched as a unit rather than via a stray subtoken. The banks are deliberately small, hand-curated lexicons---16 target cues and 24 attack cues (Appendix~\ref{sec:cuebank})---seeded from inspection of the training split and common Vietnamese expressions. Because they are seeded from training data they risk encoding train-specific patterns, and their size bounds recall on novel slang; we report the resulting coverage in Section~\ref{sec:coverage} and revisit these constraints in the Limitations.

\section{Methodology}

\subsection{Problem formulation}
Given a Vietnamese comment $x$, the model predicts $y \in \{\textsc{non-hate}, \textsc{hate}\}$. ViAmpleHate inherits AmpleHate's target-aware idea and adapts it through five changes: Vietnamese target extraction, separate target/attack cue channels, relation-bank attention, a corrected batched attention computation, and adaptive-gate injection (Figure~\ref{fig:arch}).

% FIGURE 2 — Overall architecture (TikZ, native vector).
\begin{figure*}[t]
  \centering
  \resizebox{\textwidth}{!}{%
  \begin{tikzpicture}[
      >={Latex[length=2mm]},
      font=\footnotesize,
      base/.style ={rectangle, rounded corners=2pt, draw=black!45, fill=black!6,
                    align=center, inner sep=4pt, minimum height=10mm, text width=20mm},
      novel/.style={rectangle, rounded corners=2pt, draw=orange!80!black, line width=0.9pt,
                    fill=orange!16, align=center, inner sep=4pt, minimum height=10mm, text width=22mm},
      attn/.style ={rectangle, rounded corners=2pt, draw=orange!75!black, fill=orange!8,
                    align=center, inner sep=3pt, minimum height=7mm, text width=16mm, font=\scriptsize},
      io/.style   ={rectangle, rounded corners=7pt, draw=blue!55, fill=blue!8,
                    align=center, inner sep=4pt, minimum height=10mm, text width=18mm},
      flow/.style ={->, draw=black!55, line width=0.7pt},
      lbl/.style  ={font=\scriptsize\itshape, fill=white, inner sep=1pt}
    ]
    % --- main path nodes (explicit coordinates, cm) ---
    \node[io]    (inp) at (0,0)      {Vietnamese\\comment $x$};
    \node[base]  (pre) at (2.5,0)    {Preprocessing\\{\scriptsize normalize $\cdot$ word-seg $\cdot$ tokenize}};
    \node[base]  (enc) at (5.2,1.3)  {PhoBERT\\encoder\\{\scriptsize $d{=}768$}};
    \node[novel] (det) at (5.2,-1.3) {Target \& attack\\detection\\{\scriptsize Vi-NER $+$ cues}};
    % --- relation-bank attention channels ---
    \node[attn] (re) at (8.4,1.15)  {$r_{\mathrm{exp}}$\\{\scriptsize target}};
    \node[attn] (ri) at (8.4,0.0)   {$r_{\mathrm{imp}}$\\{\scriptsize context}};
    \node[attn] (ra) at (8.4,-1.15) {$r_{\mathrm{atk}}$\\{\scriptsize attack}};
    % --- downstream ---
    \node[novel] (gate) at (11.9,0)  {Adaptive gate\\{\scriptsize $z{=}h_0{+}g\,r$}};
    \node[base]  (clf)  at (14.2,0)  {Classifier\\{\scriptsize $+$ threshold}};
    \node[io]    (out)  at (16.4,0)  {NON-HATE\\/ HATE};
    \node[font=\scriptsize] (loss) at (11.9,-2.1) {$L=L_{\mathrm{CE}}+\alpha\,L_{\mathrm{CL}}$};
    % --- relation-bank container (background) ---
    \begin{scope}[on background layer]
      \node[draw=orange!70!black, dashed, rounded corners=3pt, fill=orange!4,
            fit=(re)(ri)(ra), inner sep=5pt,
            label={[font=\scriptsize, text=orange!55!black]above:Relation-bank attention}] (relbox) {};
    \end{scope}
    % --- arrows ---
    \draw[flow] (inp) -- (pre);
    \draw[flow] (pre.east) -- (enc.west);
    \draw[flow] (pre.east) -- (det.west);
    \draw[flow] (enc.east) -- node[lbl,pos=0.65]{$h_0$} (relbox.north west);
    \draw[flow] (det.east) -- node[lbl,pos=0.65]{$T_x,A_x$} (relbox.south west);
    \draw[flow] (relbox.east) -- node[lbl]{$r$} (gate.west);
    \draw[flow] (enc.north) -- (5.2,2.3) -- node[lbl]{$h_0$ (residual)} (11.9,2.3) -- (gate.north);
    \draw[flow] (gate) -- node[lbl]{$z$} (clf);
    \draw[flow] (clf) -- (out);
    \draw[->, draw=black!40, dashed] (gate.south) -- (loss.north);
  \end{tikzpicture}%
  }
  \caption{Overall architecture of ViAmpleHate. A comment is normalized, word-segmented, and encoded by PhoBERT (global representation $h_0$); Vietnamese NER and the target-cue bank yield the target set $T_x$ while the attack-cue bank yields $A_x$. Three relation channels---explicit target ($r_{\mathrm{exp}}$), implicit context ($r_{\mathrm{imp}}$), and attack ($r_{\mathrm{atk}}$)---are fused into $r$ and injected into $h_0$ through an instance-adaptive gate ($z=h_0+g\,r$) before classification. The model is trained with weighted cross-entropy plus a contrastive term ($L=L_{\mathrm{CE}}+\alpha L_{\mathrm{CL}}$). Orange marks the components introduced by ViAmpleHate.}
  \label{fig:arch}
\end{figure*}

\subsection{Text preprocessing}
As above, comments are normalized, word-segmented, tokenized with the PhoBERT tokenizer, and truncated/padded to $\mathrm{max\_len}=256$.

\subsection{Multi-signal target and attack extraction}
The target index set combines Vietnamese NER and target-cue matching; the attack index set comes from attack-cue matching. If either is empty, the model falls back to \texttt{[CLS]} so an implicit, sentence-level representation is always available:
\begin{align}
  T_x &= M_{\mathrm{NER}}(x)\,\cup\,M_{\mathrm{target}}(x), \;\; A_x = M_{\mathrm{attack}}(x) \\
  T_x &\leftarrow \{0\}\ \text{if } T_x=\varnothing, \;\; A_x \leftarrow \{0\}\ \text{if } A_x=\varnothing \nonumber
\end{align}

\subsection{PhoBERT encoder}
\begin{equation}
  H = \mathrm{PhoBERT}(x) = [\,h_0, h_1, \dots, h_n\,], \quad h_i \in \mathbb{R}^{d},
\end{equation}
with $d=768$; $h_0$ is the \texttt{[CLS]} representation (global context), and the target/attack positions provide localized evidence.

\subsection{Relation-bank attention}
We build three relation views---explicit target, implicit context, and attack:
\begin{align}
  r_{\mathrm{exp}} &= \mathrm{HeadAttn}(h_0, H[T_x]), \nonumber\\
  r_{\mathrm{imp}} &= \mathrm{HeadAttn}(h_0, h_0), \\
  r_{\mathrm{atk}} &= \mathrm{HeadAttn}(h_0, H[A_x]), \nonumber
\end{align}
where each module over $E \in \mathbb{R}^{m\times d}$ computes
\begin{equation}
  \alpha = \mathrm{softmax}\!\Big(\tfrac{(W_q h_0)(W_k E)^\top}{\sqrt d}\Big),\;\; r = \alpha\,(W_v E).
\end{equation}
The three vectors are fused as $r = W_r[\,r_{\mathrm{exp}};r_{\mathrm{imp}};r_{\mathrm{atk}}\,] + b_r$.

\subsection{Batched attention (implementation note)}
In our re-implementation, an AmpleHate-style attention computed as $QK^\top$ over the whole batch ($Q,K\in\mathbb{R}^{B\times d}\Rightarrow QK^\top\in\mathbb{R}^{B\times B}$) can mix information across samples. We therefore use batched attention so each sample attends only to its own tokens, with $Q\in\mathbb{R}^{B\times 1\times d}$, $K\in\mathbb{R}^{B\times m\times d}$:
\begin{equation}
  \mathrm{scores} = \mathrm{bmm}(Q,K^\top)/\sqrt d \in\mathbb{R}^{B\times1\times m},
\end{equation}
with padding masked before softmax ($\mathrm{scores}_j=-\infty$ if $\mathrm{mask}_j=0$). We treat this as an implementation correction to our own port rather than a claim about the original AmpleHate, and report it as an implementation note rather than evaluating it in isolation.

\subsection{Instance-adaptive relation gate}
We replace the fixed scalar with a per-sample gate:
\begin{equation}
  g = \sigma\big(W_g[\,h_0;r\,] + b_g\big),\qquad z = h_0 + g\cdot r.
\end{equation}
With clear evidence, the model raises $g$ and relies more on the relation vector; with weak or absent cues, it lowers $g$ and leans on the global representation. The fused $z$ passes through dropout and a linear layer to produce logits.

\subsection{Training objective}
We train with $L = L_{\mathrm{CE}} + \alpha\,L_{\mathrm{CL}}$, $\alpha=0.1$. Weighted cross-entropy $L_{\mathrm{CE}}=-\sum_c w_c y_c\log\hat y_c$ handles imbalance by up-weighting HATE, with label smoothing (the class weights $w_c$, the label-smoothing value, and the margin $m$ below are reported in Appendix~\ref{sec:hparams}). The contrastive term on $z$ (with $s_{ij}=\cos(z_i,z_j)$) pulls same-class pairs together and pushes different-class pairs apart:
\begin{equation}
\begin{split}
  L_{\mathrm{CL}} = \frac{1}{N}\sum_{i\neq j}\big[&\mathbb{1}[y_i{=}y_j](1-s_{ij}) \\[-2pt]
   + &\,\mathbb{1}[y_i{\neq}y_j]\max(0, s_{ij}-m)\big].
\end{split}
\end{equation}

\subsection{Inference and threshold selection}
At inference we apply a decision threshold chosen on validation by maximizing macro-F1 and then fixed for test:
\begin{equation}
  t^{*} = \arg\max_t\ \mathrm{MacroF1}\big(y,\ \mathbb{1}[p_{\mathrm{HATE}}\ge t]\big).
\end{equation}

\subsection{Computational considerations}
The added components are lightweight relative to PhoBERT. The three HeadAttention modules and the gate operate on the pooled \texttt{[CLS]} query and a small number of cue tokens per sample, so the extra cost is a few linear projections; the corrected batched attention removes the spurious $B\times B$ interaction and is cheaper at training time. Cue matching is performed once during preprocessing.

\section{Experiments}

\subsection{Baselines}
We compare against five baselines, all binary. \textbf{TF-IDF + LR} and \textbf{TF-IDF + SVM} use sparse lexical features. \textbf{BiLSTM + fastText} \citep{fasttext} combines static embeddings with a recurrent model. \textbf{PhoBERT-CNN} stacks CNN feature extractors on PhoBERT. \textbf{AmpleHate-PhoBERT} is a port of the original with English NER, a single HeadAttention, and fixed injection $z=h_0+e\,r_{\mathrm{base}}$ ($e=1.0$). It shares the PhoBERT-base encoder. As Table~\ref{tab:arch} makes explicit, each model is run at its own intended configuration---ViAmpleHate deliberately uses a longer context window (256) and a larger effective batch, which are part of the proposed system---so the comparison is between the full proposed system and the baseline as configured, rather than an isolation of individual components.

\begin{table*}[t]
  \centering
  \small
  \begin{tabular}{p{0.20\linewidth} p{0.36\linewidth} p{0.36\linewidth}}
    \toprule
    \textbf{Component} & \textbf{Baseline AmpleHate-PhoBERT} & \textbf{ViAmpleHate-PhoBERT (proposed)} \\
    \midrule
    Encoder & PhoBERT-base & PhoBERT-base \\
    Target extraction & English NER & Vietnamese NER + target-cue bank \\
    Target coverage & $\approx$0.09\% of train (21/24{,}048) & $\approx$18.8--20.0\% via Vietnamese cues \\
    Attack signal & Not modeled & Separate attack-cue bank \\
    Attention & One HeadAttention & Three channels: target / implicit / attack \\
    Attention computation & Batch-level \texttt{matmul} (mixes samples) & Per-sample \texttt{bmm} + masking \\
    Fusion & Fixed injection $h_0+e\,r$ ($e{=}1.0$) & Relation bank + adaptive gate $h_0+g\,r$ \\
    Loss & Weighted CE & Weighted CE + contrastive ($\alpha{=}0.1$) \\
    Max length & 128 & 256 \\
    Batch & 16 & 16 $\times$ grad-accum 2 (eff.\ 32) \\
    NER at evaluation & Off ($\Rightarrow$ \texttt{[CLS]} fallback) & On, consistent across splits \\
    \bottomrule
  \end{tabular}
  \caption{Architecture and configuration comparison between the baseline and the proposed model.}
  \label{tab:arch}
\end{table*}

\subsection{Implementation details}
The encoder is \texttt{vinai/phobert-base} (768-d); Vietnamese NER is \texttt{NlpHUST/ner-vietnamese-electra-base}. Max length is 256. Learning rates are $2\mathrm{e}{-}5$ (encoder) and $5\mathrm{e}{-}5$ (head); dropout $0.1$. Effective batch is 32 (16 $\times$ grad-accum 2). We train up to eight epochs and select the best checkpoint by validation macro-F1. $\alpha=0.1$, and the threshold is selected on validation by macro-F1. Full hyperparameters are in Appendix~\ref{sec:hparams}.

\subsection{Evaluation metrics}
Our primary metrics are macro-F1 and HATE-class F1, which reflect minority-class performance far better than accuracy: under $\approx$10\% positives, a trivial majority predictor already exceeds 0.89 accuracy while detecting no hate. Macro-F1 averages per-class F1; HATE-F1 isolates the class of interest.

\subsection{Experimental design}
Beyond final scores, our experiments test whether each modification addresses a concrete baseline limitation: Vietnamese NER and target cues target low coverage; the attack-cue bank addresses missing hostility modeling; relation-bank attention separates target, context, and attack; corrected batched attention removes cross-sample leakage; the adaptive gate replaces fixed injection; and the contrastive loss improves class separation. Section~\ref{sec:coverage} examines target coverage directly.

\section{Results and Analysis}

\subsection{Main results}

Each model is reported at its own intended configuration (Section~\ref{sec:reshape}, Table~\ref{tab:arch}); the proposed-vs-baseline rows therefore compare the full proposed system against the baseline as configured. All numbers are from a single run (seed 42).

\begin{table}[t]
  \centering
  \small
  \begin{tabular}{lrrr}
    \toprule
    \textbf{Model} & \textbf{Acc.} & \textbf{Mac-F1} & \textbf{HATE-F1} \\
    \midrule
    TF-IDF + LR          & 0.8910 & 0.7393 & 0.5404 \\
    TF-IDF + SVM         & 0.9126 & 0.7131 & 0.4739 \\
    BiLSTM + fastText    & 0.8454 & 0.7072 & 0.5060 \\
    PhoBERT-CNN          & 0.8945 & 0.7571 & 0.5745 \\
    AmpleHate (baseline) & 0.9175 & 0.7792 & 0.6045 \\
    \textbf{ViAmpleHate} & \textbf{0.9205} & \textbf{0.7819} & \textbf{0.6081} \\
    \textit{$\Delta$ base} & \textit{+.0030} & \textit{+.0027} & \textit{+.0036} \\
    \bottomrule
  \end{tabular}
  \caption{Results on ViHSD (test). Best per column in bold (accuracy shown for reference). Single run (seed 42).}
  \label{tab:vihsd}
\end{table}

\begin{table}[t]
  \centering
  \small
  \begin{tabular}{lrrr}
    \toprule
    \textbf{Model} & \textbf{Acc.} & \textbf{Mac-F1} & \textbf{HATE-F1} \\
    \midrule
    TF-IDF + LR          & 0.9453 & 0.7745 & 0.5783 \\
    TF-IDF + SVM         & 0.9641 & 0.7831 & 0.5850 \\
    BiLSTM + fastText    & 0.8650 & 0.6712 & 0.4187 \\
    PhoBERT-CNN          & 0.9623 & 0.8150 & 0.6500 \\
    AmpleHate (baseline) & 0.9643 & 0.8185 & 0.6557 \\
    \textbf{ViAmpleHate} & 0.9420 & \textbf{0.8371} & \textbf{0.7065} \\
    \textit{$\Delta$ base} & \textit{--.0223} & \textit{+.0186} & \textit{+.0508} \\
    \bottomrule
  \end{tabular}
  \caption{Results on VOZ-HSD (test). Best per column in bold among models (accuracy shown for reference; the baseline has the highest accuracy). Single run (seed 42).}
  \label{tab:voz}
\end{table}

Transformer-based models clearly outperform the lexical and static-embedding baselines on the metrics that matter (Tables~\ref{tab:vihsd},~\ref{tab:voz}), confirming that contextual representations are important when hate depends on context, informal phrasing, and the interaction between a mention and a hostile predicate. Within the transformer family, the AmpleHate baseline improves over a plain PhoBERT classifier through target-aware attention, but its benefit is capped because English NER rarely detects valid Vietnamese targets, so most inputs revert to \texttt{[CLS]}. We compare only against models we trained under a common protocol; a comparison to published ViHSD results is left to future work.

ViAmpleHate improves macro-F1 and HATE-F1 on VOZ-HSD by a clear margin (HATE-F1 $+0.0508$). On ViHSD the differences are small ($+0.0027$ macro-F1, $+0.0036$ HATE-F1); since all reported numbers come from a single run (seed 42) and we do not run a significance test, we treat the ViHSD margin as preliminary rather than established. Accuracy \emph{drops} on VOZ-HSD while macro-F1 and HATE-F1 rise, illustrating why accuracy is the wrong headline metric under imbalance.

\subsection{Effect of Vietnamese target extraction}
\label{sec:coverage}
A likely contributor to these gains is target coverage, which we define as the fraction of comments for which $T_x \neq \{0\}$ (Vietnamese NER or a target cue fires). Under English NER, only $\approx$0.09\% of training comments contain a detected target (21/24{,}048), so the baseline's target-aware pathway is almost never active. With Vietnamese NER and the target-cue bank, target coverage rises to 20.0\%/18.8\% on the first 500 ViHSD train/validation comments and to 45.2\%/43.0\% on VOZ-HSD using the same 16-cue bank; attack-cue coverage is lower---5.2\%/4.2\% (ViHSD) and 7.6\%/6.0\% (VOZ). Notably, the dataset with far higher coverage (VOZ-HSD) is also where ViAmpleHate gains most (HATE-F1 $+0.0508$ vs.\ $+0.0036$), a suggestive cross-dataset correlation. We stress, however, that coverage is an \emph{input} statistic: higher coverage gives relation-bank attention more to attend to but does not by itself prove that firing improves individual decisions.

% FIGURE 3 — Training curves: training_curves_viamplehate.png
\begin{figure}[t]
  \centering
  \includegraphics[width=\columnwidth]{example-image-a}
  \caption{Training/validation curves of ViAmpleHate on ViHSD (best epoch 4, val macro-F1 $=0.7852$). \textbf{[USE PNG]} \texttt{training\_curves\_viamplehate.png}.}
  \label{fig:curves}
\end{figure}

\subsection{Qualitative analysis}
To make the error modes concrete, Table~\ref{tab:qual} shows illustrative comment types (constructed examples that mirror the observed categories; replace with real anonymized test instances for camera-ready). The pattern is that ViAmpleHate helps most when an explicit target co-occurs with an attack predicate, and least for implicit hate that names no target and uses no overt attack term.

\begin{table}[t]
  \centering
  \small
  \begin{tabular}{p{0.34\linewidth}ccc}
    \toprule
    \textbf{Comment type} & \textbf{Tgt} & \textbf{Atk} & \textbf{Gold} \\
    \midrule
    Group label + hostile predicate     & yes & yes & HATE \\
    Profanity, no group target           & no  & yes & NON \\
    Sarcastic/implicit, no overt cue     & no  & no  & HATE \\
    Group label in neutral context       & yes & no  & NON \\
    \bottomrule
  \end{tabular}
  \caption{Illustrative cases (constructed; replace with real examples). Tgt/Atk = target/attack cue present.}
  \label{tab:qual}
\end{table}

% FIGURE 5 — Confusion matrices ViHSD: baseline vs proposed.
\begin{figure*}[t]
  \centering
  \includegraphics[width=0.45\linewidth]{example-image-a}\hfill
  \includegraphics[width=0.45\linewidth]{example-image-b}
  \caption{Confusion matrices on ViHSD: baseline AmpleHate (left) vs ViAmpleHate (right). \textbf{[USE PNGs]} \texttt{confusion\_matrix\_amplehate.png} and \texttt{confusion\_matrix\_viamplehate.png}; note the reduction in HATE false positives.}
  \label{fig:cm}
\end{figure*}

\begin{table}[t]
  \centering
  \small
  \begin{tabular}{llrrr}
    \toprule
    \textbf{Data} & \textbf{Class} & \textbf{P} & \textbf{R} & \textbf{F1} \\
    \midrule
    ViHSD & NON-HATE & 0.9541 & 0.9574 & 0.9558 \\
    ViHSD & HATE     & 0.6177 & 0.5988 & 0.6081 \\
    VOZ   & NON-HATE & 0.9638 & 0.9719 & 0.9678 \\
    VOZ   & HATE     & 0.7347 & 0.6803 & 0.7065 \\
    \bottomrule
  \end{tabular}
  \caption{Per-class results of ViAmpleHate (test set, single run).}
  \label{tab:perclass}
\end{table}

\subsection{Discussion}
Three observations stand out. First, on ViHSD the change is a precision/recall reweighting rather than a uniform gain: HATE precision rises ($0.5972\!\to\!0.6177$) while recall \emph{falls} ($0.6119\!\to\!0.5988$). Higher precision is desirable when false accusations of hate are costly, but lower recall means more hate is missed; which trade-off is appropriate depends on the deployment. We report results at a single validation-tuned threshold and leave a full precision--recall analysis to future work. On VOZ-HSD, where coverage is higher, HATE recall is also higher (0.6803; Table~\ref{tab:perclass}). Second, the larger VOZ-HSD gains alongside an accuracy drop indicate that the model trades majority-class accuracy for minority-class quality, which is the trade-off practitioners typically want when the minority class is the target of interest. Third, the relationship between cue coverage and gains suggests a possible inexpensive path to further improvement---expanding and refining the cue banks---though this remains to be verified (Section~\ref{sec:coverage}).

\section{Error Analysis}
We group the remaining errors into seven recurring categories. \textbf{(1) Offensive but not hate:} insults or profanity aimed at an individual rather than a protected group; over-weighting attack cues yields false positives. \textbf{(2) Implicit hate:} comments with no slur, named entity, or overt attack predicate, expressing hostility through sarcasm, stereotypes, or shared social context; cue extraction finds little to attend to. \textbf{(3) Ambiguous target reference:} pronouns and group nouns also occur in neutral or humorous comments, so a target cue without hostile context can over-predict HATE. \textbf{(4) Incomplete cue coverage:} fast-changing slang, spelling variants, abbreviations, and creative or censored profanity that a fixed cue bank cannot cover, causing false negatives. \textbf{(5) Tokenization and span-alignment errors:} NER spans, word segmentation, and PhoBERT subword tokenization may misalign---especially for multi-word expressions---so a cue can match the wrong position. \textbf{(6) Threshold sensitivity:} a validation-tuned threshold may not transfer under distribution shift. \textbf{(7) Class imbalance:} with few positives, minority-class errors persist despite weighted loss and threshold tuning. These point to expanding the cue banks via data-driven mining plus manual validation, logging per-instance predictions (with gate value, cue coverage, and confidence) for systematic false-positive/negative analysis, adding span supervision or sarcasm detection, and applying probability calibration.

\section{Conclusion and Future Work}
ViAmpleHate adapts AmpleHate to Vietnamese hate speech detection through Vietnamese target extraction, separate target/attack cue channels, relation-bank attention, and adaptive relation injection (plus an implementation-level batched-attention correction), trained with weighted cross-entropy plus a contrastive loss. On ViHSD and VOZ-HSD it improves macro-F1 and HATE-F1 over the AmpleHate baseline and over the TF-IDF, BiLSTM, and PhoBERT-CNN baselines, with the clearest gains on VOZ-HSD's minority class; the ViHSD improvements are small and remain to be confirmed under a controlled, multi-seed comparison. Our analysis relates these gains to target coverage, which the Vietnamese cue banks raise by roughly two orders of magnitude over English NER---a correlational link that we have not yet shown to be causal. The broader lesson is that target-awareness is useful only when adapted to the target language's surface forms, where hate targets are often informal group references rather than named entities. Future work will expand and automatically mine the cue banks with manual validation; analyse errors at the per-instance level using gate value and confidence; add span supervision, sarcasm detection, and user/discourse context; extend to multi-label and multi-target settings in the spirit of ViTHSD; and apply probability calibration for stable cross-domain deployment.

\section*{Limitations}
The target and attack cue banks are partly hand-built and cannot cover the full, fast-changing space of Vietnamese slang, spelling variants, and creative profanity, which bounds recall on implicit or novel hate. Performance depends on upstream Vietnamese NER, word segmentation, and tokenization, whose misalignment can misplace cues. The decision threshold is tuned on validation and may not be optimal under distribution shift, and reporting at a single threshold does not expose the full precision--recall trade-off. All results come from a single run (seed 42); we do not report multi-seed variance or significance tests, and we do not ablate the individual components, so component-level contributions remain design hypotheses and the small ViHSD margin is not established. We compare against our own re-implementations rather than published ViHSD results. Finally, the model addresses a binary NON-HATE/HATE formulation and does not yet handle multi-target or graded hatred.

\section*{Ethics Statement}
Hate speech detection is dual-use: the same models that support content moderation can, if mis-deployed, suppress legitimate speech or disproportionately flag particular communities. Our datasets are existing, publicly described Vietnamese resources used under their intended research terms; we do not collect new user data. Because annotations and cue banks reflect particular annotator and author judgments, the system may carry biases against specific dialects, regions, or groups, and its outputs should be treated as decision support for human moderators rather than automated enforcement. We avoid reproducing real slurs in the paper, using constructed or masked examples. We encourage calibration, human-in-the-loop review, and ongoing bias auditing in any deployment.

% \section*{Acknowledgments}

\bibliography{custom}

\appendix

\section{Hyperparameters}
\label{sec:hparams}
\begin{table}[h]
  \centering
  \small
  \setlength{\tabcolsep}{3pt} 
  \begin{tabular}{@{}>{\raggedright\arraybackslash}p{0.50\columnwidth}>{\raggedright\arraybackslash}p{0.47\columnwidth}@{}} 
    \toprule
    \textbf{Parameter} & \textbf{Value} \\
    \midrule
    Encoder & \texttt{vinai/phobert-base} (768-d) \\
    Vietnamese NER & \texttt{NlpHUST/\allowbreak ner-vietnamese-\allowbreak electra-base} \\
    max\_len & 256 \\
    LR (encoder / head) & 2e-5 / 5e-5 \\
    Dropout & 0.1 \\
    Effective batch & 32 (16 $\times$ grad-accum 2) \\
    Epochs & up to 8 (best by val macro-F1) \\
    $\alpha$ (contrastive) & 0.1 \\
    Class weights $w_c$ & inverse freq.\ $N/(C\,n_c)$ \\
    Contrastive margin $m$ & 0.5 \\
    Label smoothing & 0.05 \\
    Seed / runs & 42 / single run \\
    ViHSD best epoch / val F1 / thr. & 4 / 0.7852 / 0.43 \\
    \bottomrule
  \end{tabular}
  \caption{Full hyperparameters.}
  \label{tab:hparams}
\end{table}

\section{Cue Bank Examples}
\label{sec:cuebank}
The full banks---16 target cues and 24 attack cues---are released with the code and listed in full in the accompanying report; we give ASCII-safe examples here. \textbf{Target cues} are referential (not hateful by themselves): informal group/person references, regional and demographic labels, group pronouns. \textbf{Attack cues} are hostile predicates expressing contempt, dehumanization, or parasitism (e.g., \texttt{khinh}, \texttt{an\_bam}). Target and attack cues are kept in separate banks (Section~\ref{sec:cue}).

\section{Reproducibility}
\label{sec:repro}
All PhoBERT-based models share \texttt{vinai/phobert-base}; only target/relation modeling, attention computation, fusion, and loss differ between the baseline and ViAmpleHate. Best checkpoints are selected by validation macro-F1, and the HATE threshold is selected on validation and fixed for test. Cue matching is performed once during preprocessing using full token-sequence matching against the PhoBERT token stream.

\end{document}
